\documentclass[11pt]{article}
% Shared print layout. The web's provisional design and MathJax are independent.
\usepackage[utf8]{inputenc}
\usepackage[T1]{fontenc}
\usepackage[spanish,es-nodecimaldot]{babel}
\usepackage{lmodern}
\usepackage{amsmath,amssymb}
\usepackage{graphicx}
\usepackage{geometry}
\geometry{a4paper,left=20mm,right=20mm,top=25mm,bottom=22mm,headheight=16pt,headsep=8mm}
\usepackage{fancyhdr,lastpage}
\usepackage{enumitem}
\usepackage{xcolor}
\usepackage{pgfplots}
\pgfplotsset{compat=1.18}
\usepgfplotslibrary{fillbetween}
\usepackage[hidelinks,pdfencoding=auto]{hyperref}
\hypersetup{pdflang=es-ES,pdfauthor={},pdfsubject={Matemáticas Aplicadas a las Ciencias Sociales II}}
\pagestyle{fancy}
\fancyhf{}
\fancyhead[L]{\small Matemáticas Aplicadas a las Ciencias Sociales II}
\fancyhead[R]{\small 2.\textsuperscript{o} Bachillerato}
\fancyfoot[L]{\footnotesize Matrices II | Edición revisada 1.0}
\fancyfoot[R]{\footnotesize Página \thepage\ de \pageref{LastPage}}
\renewcommand{\headrulewidth}{0.4pt}
\setlength{\parindent}{0pt}
\setlength{\parskip}{0.5em}
\setlist{itemsep=0.3em,topsep=0.3em}
\newcommand{\apartado}[1]{\subsection*{#1}}
\pgfplotsset{
  revision/.style={
    width=0.94\linewidth,height=7.5cm,
    axis lines=middle,grid=major,grid style={gray!18},
    tick label style={font=\small},label style={font=\small},
    xlabel={$x$},ylabel={$y$},samples=160,
    legend style={at={(0.5,-0.16)},anchor=north,draw=none,font=\small},
    clip=true
  }
}

\hypersetup{pdftitle={Matrices II - Soluciones revisadas}}
\begin{document}
\section*{Matrices II: soluciones}
Solucionario elaborado en esta revisión a partir de F015. El ZIP original no incluía soluciones.
\section*{Ejercicio 1(a): determinante e inversa}
Para $A=\begin{pmatrix}1&1&3\\1&2&-1\\1&-1&-1\end{pmatrix}$, desarrollamos por la primera fila:
\[
\det A=1(-2-1)-1(-1+1)+3(-1-2)=-12\neq0.
\]
Por tanto $A$ tiene inversa. Su matriz de cofactores es
\[
K=\begin{pmatrix}-3&0&-3\\-2&-4&2\\-7&4&1\end{pmatrix}.
\]
Transponemos para obtener la adjunta y dividimos por el determinante:
\[
A^{-1}=\frac1{-12}K^t
=\boxed{\begin{pmatrix}1/4&1/6&7/12\\0&1/3&-1/3\\1/4&-1/6&-1/12\end{pmatrix}}.
\]
\textbf{Comprobación:} $AA^{-1}=A^{-1}A=I_3$.
\section*{Ejercicio 1(b): resolver $AX=B$}
$B$ es $3\times1$, por lo que $X$ debe ser $3\times1$. Multiplicamos por $A^{-1}$ a la izquierda:
\[
X=A^{-1}B=
\begin{pmatrix}1/4&1/6&7/12\\0&1/3&-1/3\\1/4&-1/6&-1/12\end{pmatrix}
\begin{pmatrix}2\\3\\0\end{pmatrix}
=\boxed{\begin{pmatrix}1\\1\\0\end{pmatrix}}.
\]
\textbf{Comprobación:}
\[
A\begin{pmatrix}1\\1\\0\end{pmatrix}
=\begin{pmatrix}1+1\\1+2\\1-1\end{pmatrix}
=\begin{pmatrix}2\\3\\0\end{pmatrix}=B.
\]
La solución existe y es única por la invertibilidad de $A$.
\clearpage
\section*{Ejercicio 2: operaciones y dimensiones}
Se muestran las cinco respuestas para consulta. En el examen se deben elegir \textbf{solo cuatro}; la puntuación total del ejercicio es 2,5 puntos.

Aquí $A,B$ son $2\times2$, $C$ es $2\times1$ y $D$ es $1\times2$. Se utilizan las matrices propias del ejercicio 2.
\subsection*{1. $AC+BD^t$}
\[
AC=\begin{pmatrix}-1\\1\end{pmatrix},\qquad
BD^t=\begin{pmatrix}-2\\5\end{pmatrix}.
\]
Los dos productos son $2\times1$:
\[\boxed{AC+BD^t=\begin{pmatrix}-3\\6\end{pmatrix}}.\]
\subsection*{2. $A+BC$}
$BC$ existe y es $2\times1$. Como $A$ es $2\times2$, la suma \textbf{no está definida}: no coinciden sus dimensiones.
\subsection*{3. $B^2+CD$}
\[
B^2=\begin{pmatrix}4&0\\-4&4\end{pmatrix},\qquad
CD=\begin{pmatrix}0&0\\-1&2\end{pmatrix}.
\]
Ambos son $2\times2$, y
\[\boxed{B^2+CD=\begin{pmatrix}4&0\\-5&6\end{pmatrix}}.\]
\subsection*{4. $|5A|$}
Para una matriz de orden 2, $|5A|=5^2|A|$. Aquí $|A|=1$:
\[\boxed{|5A|=25}.\]
\subsection*{5. $|C|$}
$C$ es $2\times1$. No es cuadrada y su determinante \textbf{no está definido}; no se le asigna el valor cero.
\clearpage
\section*{Ejercicio 3: igualdad con parámetros}
Multiplicamos ambos lados, prestando atención a que la segunda coordenada de la columna izquierda es $-b$:
\[
\begin{pmatrix}2&-1\\3&-1\end{pmatrix}\begin{pmatrix}a\\-b\end{pmatrix}
=\begin{pmatrix}2a+b\\3a+b\end{pmatrix},
\]
\[
\begin{pmatrix}1&2\\b&-1\end{pmatrix}\begin{pmatrix}3\\0\end{pmatrix}
=\begin{pmatrix}3\\3b\end{pmatrix}.
\]
Dos columnas son iguales cuando coinciden sus coordenadas:
\[
\begin{cases}2a+b=3,\\3a+b=3b.\end{cases}
\quad\Longleftrightarrow\quad
\begin{cases}2a+b=3,\\3a-2b=0.\end{cases}
\]
De la segunda ecuación, $b=3a/2$. Sustituyendo en la primera,
\[2a+\frac32a=3\quad\Longrightarrow\quad\frac72a=3\quad\Longrightarrow\quad a=\frac67.\]
Así,
\[\boxed{a=\frac67,\qquad b=\frac97}.\]
\textbf{Comprobación en la igualdad original:}
\[
\begin{pmatrix}2a+b\\3a+b\end{pmatrix}
=\begin{pmatrix}3\\27/7\end{pmatrix}
=\begin{pmatrix}3\\3b\end{pmatrix}.
\]
El sistema de parámetros tiene determinante $2(-2)-3\cdot1=-7\neq0$; la pareja es única.
\end{document}
